\documentclass[11pt,a4paper]{article}
\usepackage[utf8]{inputenc}
\usepackage[T1]{fontenc}
\usepackage[french]{babel}
\usepackage[margin=2.4cm]{geometry}
\usepackage{amsmath,amssymb}
\usepackage{booktabs}
\usepackage{xcolor}
\usepackage[colorlinks=true,linkcolor=blue!50!black,citecolor=blue!50!black]{hyperref}
\usepackage{tcolorbox}
\tcbuselibrary{skins,breakable}
\usepackage{titlesec}
\usepackage{enumitem}
\usepackage{parskip}
\setlength{\parskip}{0.4em}

\definecolor{bleuprincip}{HTML}{1F4E78}
\definecolor{bleuclair}{HTML}{2E74B5}
\definecolor{orangealerte}{HTML}{C55A11}
\definecolor{vertok}{HTML}{548235}

\titleformat{\section}{\Large\bfseries\color{bleuprincip}}{\thesection}{0.6em}{}
\titleformat{\subsection}{\large\bfseries\color{bleuclair}}{\thesubsection}{0.6em}{}

\newtcolorbox{intuition}[1][]{colback=bleuclair!7,colframe=bleuclair,boxrule=0.5pt,arc=2pt,
  left=6pt,right=6pt,top=5pt,bottom=5pt,breakable,title={\small\bfseries Intuition},#1}
\newtcolorbox{definitionbox}[1][]{colback=gray!7,colframe=black!60,boxrule=0.4pt,arc=1pt,
  left=6pt,right=6pt,top=5pt,bottom=5pt,breakable,#1}
\newtcolorbox{resultatbox}[1][]{colback=vertok!8,colframe=vertok,boxrule=1pt,arc=2pt,
  left=7pt,right=7pt,top=5pt,bottom=5pt,breakable,#1}
\newtcolorbox{alertebox}[1][]{colback=orangealerte!7,colframe=orangealerte,boxrule=0.6pt,arc=2pt,
  left=7pt,right=7pt,top=5pt,bottom=5pt,breakable,title={\small\bfseries Résultat notable},#1}

\newcommand{\cit}[1]{\textsuperscript{\hyperlink{ref#1}{[#1]}}}

\title{\vspace{-1.3cm}\textcolor{bleuprincip}{\Huge\bfseries GDPNow-Maroc}\\[0.3cm]
\textcolor{black}{\Large Méthode 3 --- Factor-Augmented Rolling AR}\\[0.15cm]
\textcolor{gray}{\large Réplique du modèle 2 de Higgins (2014), Table 5}}
\author{}
\date{\today}

\begin{document}
\maketitle
\thispagestyle{empty}

\begin{abstract}
\noindent Ce document construit le troisième modèle de comparaison : un AR augmenté du facteur commun de chaque branche (Étape 2b, \emph{Method\_1\_Higgins}), avec $(q,r)$ choisis par AIC --- l'équivalent du modèle 2 de Higgins\cit{1} (« Rolling Factor Augmented AR(2) »). Sur la même fenêtre de test que les Méthodes 1 et 2 (2021T2--2026T1), il obtient un RMSFE de 12,79 --- meilleur que l'AR pur (Méthode 2, 14,55) mais toujours nettement en retrait du modèle complet (Méthode 1, 5,43). Le même classement relatif que Higgins observe dans son propre article.
\end{abstract}

\tableofcontents
\vspace{0.5cm}

% ============================================================================
\section{Méthode}
% ============================================================================

\begin{equation}
\Delta\log(VA_t) = \alpha + \sum_{k=1}^{q}\gamma_k \Delta\log(VA_{t-k}) + \sum_{j=0}^{r}\beta_j F_{t-j} + \varepsilon_t,
\qquad (q^\star,r^\star) = \operatorname*{arg\,min}_{q,r\in\{0,\dots,8\}} \text{AIC}(q,r)
\end{equation}

\begin{intuition}
Higgins fixe $q=2$ pour ce modèle de comparaison. Comme pour la Méthode 2, ce choix arbitraire n'est pas reproduit --- $(q,r)$ sont recherchés par AIC, indépendamment pour chacune des branches concernées, cohérent avec toute la rigueur déjà appliquée dans ce travail.
\end{intuition}

\subsection{Couverture --- seulement 11 branches sur 16 en Factor-AR}

\begin{definitionbox}
Le facteur commun n'existe que là où au moins 2 indicateurs ont été retenus en Phase 2 --- 12 branches (Étape 2b, Modèle 1). Les 4 branches sans aucun indicateur (Services aux entreprises, Administration publique, Éducation-santé, Autres services) restent en \textbf{AR pur} (repris tel quel de la Méthode 2, jamais réestimé).
\end{definitionbox}

\begin{alertebox}
\textbf{Agriculture, bien que disposant d'un facteur, en est exclue elle aussi} : son facteur (précipitations/température) ne couvre que 2020--2026 --- une fois croisé avec le train (2010--2021), il ne reste que 4--5 trimestres exploitables, très en dessous du minimum requis pour une recherche $(q,r)$ fiable sur une grille $9\times9$. Repli sur l'AR pur (Méthode 2), documenté honnêtement plutôt que forcé.
\end{alertebox}

\textbf{Bilan final} : 11 branches en Factor-AR véritable, 5 en AR pur (Agriculture + les 4 sans indicateur).

% ============================================================================
\section{Résultats}
% ============================================================================

\subsection{Les 11 couples $(q^\star,r^\star)$ retenus}

\begin{center}
\small
\begin{tabular}{lrrrr}
\toprule
\textbf{Branche} & \textbf{$q^\star$} & \textbf{$r^\star$} & \textbf{$R^2$} & \textbf{$n$} \\
\midrule
Hébergement-restauration & 5 & 4 & 0,752 & 45 \\
Industrie de transformation & 7 & 1 & 0,798 & 37 \\
Transports & 7 & 3 & 0,611 & 40 \\
Industrie d'extraction & 8 & 4 & 0,591 & 45 \\
Construction & 4 & 3 & 0,504 & 45 \\
Pêche & 1 & 0 & 0,398 & 43 \\
Information-communication & 3 & 5 & 0,309 & 45 \\
Immobilier & 4 & 1 & 0,247 & 40 \\
Électricité, gaz, eau & 1 & 0 & 0,203 & 43 \\
Finances et assurances & 0 & 0 & 0,232 & 44 \\
Commerce & 1 & 2 & 0,164 & 38 \\
\bottomrule
\end{tabular}
\end{center}

\subsection{En-échantillon vs hors-échantillon}

\begin{resultatbox}
\begin{center}
\begin{tabular}{lrr}
\toprule
& \textbf{En-échantillon (train)} & \textbf{Hors-échantillon (test)} \\
\midrule
Corrélation & 0,543 & $-0{,}224$ \\
$n$ & 44 & 20 \\
\bottomrule
\end{tabular}
\end{center}
\end{resultatbox}

% ============================================================================
\section{Comparaison complète --- les trois méthodes et Higgins}
% ============================================================================

\begin{center}
\small
\begin{tabular}{lrrl}
\toprule
\textbf{Modèle} & \textbf{RMSFE} & \textbf{Corrélation} & \textbf{Contexte} \\
\midrule
Higgins --- Rolling AR(2) & 2,43 & --- & USA, $q=2$ fixe \\
Higgins --- Factor Augmented AR(2) & 1,95 & --- & USA, $q=2$ fixe \\
Higgins --- GDPNow complet & \textbf{1,15} & --- & USA \\
\midrule
Méthode 2 (nous) --- Rolling AR & 14,55 & $-0{,}060$ & Maroc, $p^\star$ par AIC \\
\textbf{Méthode 3 (nous) --- Factor-Augmented AR} & \textbf{12,79} & $-0{,}224$ & Maroc, $(q^\star,r^\star)$ par AIC \\
Méthode 1 (nous) --- modèle complet & \textbf{5,43} & 0,292 & Maroc \\
\bottomrule
\end{tabular}
\end{center}

\begin{resultatbox}
\textbf{Le même classement relatif que Higgins se confirme} : Factor-Augmented AR bat Rolling AR simple, chez lui (1,95 contre 2,43, un gain de 20\%) comme chez nous (12,79 contre 14,55, un gain de 12\%) --- \textbf{mais dans les deux cas, ni l'un ni l'autre n'approche le modèle complet}. Le facteur apporte un vrai gain, réel mais modeste, jamais suffisant à lui seul.
\end{resultatbox}

\begin{alertebox}
\textbf{Une nuance à souligner} : le gain en RMSFE (12,79 contre 14,55) s'accompagne d'une corrélation \emph{plus négative} ($-0{,}224$ contre $-0{,}060$) --- RMSFE et corrélation mesurent des choses différentes (l'ampleur des erreurs contre le sens du co-mouvement). Le facteur réduit l'amplitude des erreurs les plus extrêmes, sans pour autant améliorer la capacité du modèle à prédire correctement la direction du mouvement.
\end{alertebox}

\section*{Références}
\addcontentsline{toc}{section}{Références}
\begin{enumerate}[leftmargin=1.6em, label=\textbf{[\arabic*]}]
\item \hypertarget{ref1}{} Higgins, P. (2014). \textit{GDPNow: A Model for GDP ``Nowcasting''}. Federal Reserve Bank of Atlanta, Working Paper 2014-7, Table 5.
\end{enumerate}

\end{document}
